% ══════════════════════════════════════════════════════════════
%  SECTION III: THE FIRST STAGE — WHAT HAPPENED TO MATCH TIME
% ══════════════════════════════════════════════════════════════

\section{The First Stage: What Happened to Match Time}
\label{sec:first_stage}

Law~7 of the Laws of the Game has required referees to add time lost to
interruptions since 1886, yet for decades the rule was systematically
under-enforced: referees added three to five minutes regardless of
actual dead time, which typically exceeded 28~minutes per match.
In November 2022, Pierluigi Collina, chairman of FIFA's Referees
Committee, instructed World Cup officials to restore the \textit{actual}
time lost---producing matches exceeding 100~minutes before a global
audience of five billion viewers.
IFAB endorsed the practice at its March 2023 Annual General Meeting,
using ``should'' rather than ``must,'' preserving each federation's
discretion and creating the institutional heterogeneity that drives
our empirical design.

After the World Cup, matches across Europe got longer---but not
uniformly, and not everywhere. The enforcement gradient that emerged
tracks institutional proximity to FIFA's refereeing apparatus with
remarkable fidelity, revealing how a global governing body's signal
travels through the decentralised architecture of European football.

We study 15~European leagues spanning the five major top-flight
competitions (the ``Big~Five'': England's Premier League, Spain's
La~Liga, Germany's Bundesliga, Italy's Serie~A, and France's
Ligue~1), their associated second divisions where data are available,
and several smaller top-flight leagues.


% ──────────────────────────────────────────────────────────────
\subsection{What Changed, and How Much}
\label{sec:first_stage:magnitude}
% ──────────────────────────────────────────────────────────────

The headline is unambiguous. Across the five treated leagues, mean
second-half stoppage time rose from 4.99~minutes ($N = 10{,}763$) to
6.56~minutes ($N = 4{,}596$), an increase of 1.57~minutes. The entire
distribution shifted rightward: a Kolmogorov-Smirnov test rejects the
null that the pre- and post-directive samples share a common distribution
($D = 0.332$, $p < 10^{-10}$). The probability of exceeding seven
minutes of added time rose from 14.8\% to 39.2\%; the probability of
exceeding ten minutes---once a rarity that generated tabloid
headlines---rose from 1.8\% to 7.8\% (Figure~\ref{fig:distributions}).

The standard deviation widened from 1.79 to 2.24~minutes. The new norm
did not replace one rigid convention with another. The fourth official's
added-time board, which once displayed a narrow range, now shows wider
variation, reflecting finer differentiation among match circumstances.
The directive asked referees to \textit{measure} rather than
\textit{estimate}, and measurement produces more variance than
convention.

Table~\ref{tab:enforcement} reports the league-by-league breakdown,
ordered by Cohen's~$d$.

\begin{table}[t]
\centering
\caption{Enforcement Intensity by League}
\label{tab:enforcement}
\small
\begin{threeparttable}
\begin{tabular}{llccccl}
\toprule
League & Code & Pre Mean & Post Mean & $\Delta$ & Cohen's $d$ & Intensity \\
\midrule
Premier League$^\dagger$ & E0 & 5.39 & 7.50 & +2.11 & 1.13 & High \\
Ligue 1$^\dagger$ & F1 & 4.71 & 6.50 & +1.79 & 1.05 & High \\
Bundesliga$^\dagger$ & D1 & 4.34 & 6.12 & +1.78 & 0.92 & High \\
Scottish Premiership & SC0 & 2.83 & 4.30 & +1.47 & 0.83 & High \\
La Liga$^\dagger$ & SP1 & 5.18 & 6.64 & +1.46 & 0.70 & High \\
Belgian Pro League & B1 & 3.09 & 4.40 & +1.31 & 0.66 & High \\
2. Bundesliga & D2 & 3.95 & 4.98 & +1.04 & 0.59 & High \\
Championship & E1 & 5.67 & 6.70 & +1.03 & 0.59 & High \\
Primeira Liga & P1 & 5.97 & 6.94 & +0.97 & 0.42 & Moderate \\
Ligue 2 & F2 & 2.81 & 3.76 & +0.95 & 0.56 & Moderate \\
Greek Super League & G1 & 3.58 & 4.41 & +0.82 & 0.38 & Moderate \\
Serie A$^\dagger$ & I1 & 5.18 & 5.91 & +0.73 & 0.40 & Moderate \\
Eredivisie & N1 & 5.24 & 5.71 & +0.47 & 0.23 & Low \\
Serie B & I2 & 4.05 & 3.97 & -0.08 & -0.04 & Low \\
\bottomrule
\end{tabular}
\begin{tablenotes}\small
\item \textit{Notes:} Pre = seasons before 2023-24. Post = 2023-24 onward. $\Delta$ = Post $-$ Pre mean second-half stoppage time (minutes). Cohen's $d$ = standardized effect size. Intensity classified as High ($\Delta > 1.0$), Moderate ($0.5 < \Delta \leq 1.0$), Low ($\Delta \leq 0.5$). $^\dagger$ denotes Big Five league.
\end{tablenotes}
\end{threeparttable}
\end{table}


The gradient is steep: the Premier League
shifted by 2.11~minutes ($d = 1.13$), more than a full standard
deviation above the pre-directive mean; Serie~B showed no measurable
change ($\Delta = -0.08$, $d = -0.04$). Between these poles, thirteen
leagues arranged themselves along a continuum that tracks institutional
proximity to FIFA with remarkable precision
(Figures~\ref{fig:enforcement_gradient} and~\ref{fig:time_series}).


% ──────────────────────────────────────────────────────────────
\subsection{How the Signal Traveled}
\label{sec:first_stage:channels}
% ──────────────────────────────────────────────────────────────

The signal originated with Pierluigi Collina. As chairman of FIFA's
Referees Committee, Collina is simultaneously the sport's most
recognizable official and its most powerful referee administrator.
When he told the 36 match officials assembled for Qatar that every
lost second must be restored, the instruction carried the weight of
both formal authority and unmatched personal credibility.
England--Iran ran to 117~minutes. The World Cup's cumulative viewership
exceeded five billion, and every viewer saw what the new standard
looked like: fourth officials holding up boards reading ``$+$10,''
``$+$12,'' ``$+$14.''

Those referees returned to domestic duty in January 2023.
Approximately 35 were members of UEFA's Elite Panel; their visible
adoption of the new norm in league matches broadcast to millions
created a public expectation that the rest of each country's referee
corps would follow. The formal institutional step came at IFAB's
Annual General Meeting in March 2023, where the board endorsed the
practice---using ``should,'' not ``must,'' and thereby preserving each
federation's discretion over implementation.

From IFAB, the signal passed through three visible channels.
\textit{Direct referee exposure}: officials who worked the World Cup
or regularly officiate FIFA fixtures received the instruction firsthand
and carried it home. \textit{Federation guidance}: pre-season
briefings, updated assessor protocols, and formal memoranda translated
IFAB's recommendation into operational standards. The assessor's
report---the post-match evaluation that grades referee performance and
influences future assignments---is the critical lever; when assessors
began flagging insufficient stoppage time as a deficiency, the
incentive facing every referee in the pool shifted overnight.
\textit{Media pressure}: in leagues with large broadcast audiences,
any referee who conspicuously failed to add adequate time faced
immediate public criticism. The pundit's question---``Why only four
minutes when there were three VAR checks?''---became a form of
decentralized enforcement requiring no monitoring apparatus and no
sanctions.


% ──────────────────────────────────────────────────────────────
\subsection{Within-Country Pairs}
\label{sec:first_stage:pairs}
% ──────────────────────────────────────────────────────────────

The most revealing comparisons are within countries, where top-flight
and second-division leagues share a national federation but differ in
their connection to FIFA's referee infrastructure.

\paragraph{England: PGMOL versus the EFL.}
The Premier League added 2.11~minutes ($d = 1.13$); the Championship
added 1.03~minutes ($d = 0.59$). The gap of more than a full minute
per match reflects the institutional divide between two referee
organizations operating under the same Football Association. The
Professional Game Match Officials Limited (PGMOL) manages Premier League
referees as full-time professionals---a status they have held since 2001,
when the FA abandoned the pretense that officials overseeing matches
worth hundreds of millions of pounds could do so as a weekend avocation.
PGMOL referees receive centralized briefings, structured assessor
reviews, and regular international assignments that connect them directly
to FIFA. The most experienced were at the World Cup or serve on UEFA's
Elite Panel; they returned to domestic duty already calibrated to the new
standard. Championship referees, managed by the EFL, sit one rung below:
not full-time professionals, but absorbing PGMOL norms through the
development pipeline, as the most promising are groomed for promotion to
the top-flight list. The Championship's substantial shift confirms the
norm penetrated beyond the treated tier; the 1.08-minute gap confirms
the signal attenuated at the organizational boundary between the two
bodies.

\paragraph{Germany: the DFB pipeline.}
The Bundesliga shifted by 1.78~minutes ($d = 0.92$); the
2.~Bundesliga by 1.04~minutes ($d = 0.59$). Germany's referee
governance is more centralized than England's: the Deutscher
Fu{\ss}ball-Bund (DFB) oversees both competitions through a unified
referee department, and promising second-division referees are developed
explicitly for Bundesliga promotion through a structured pathway. The
within-Germany gap (0.74~minutes) is smaller than the within-England
gap (1.08~minutes), consistent with the DFB's unified pipeline
transmitting the enforcement signal more efficiently to its second
division than the FA's bifurcated structure manages.

\paragraph{France: shared officials.}
Ligue~1 added 1.79~minutes ($d = 1.05$); Ligue~2 added 0.95~minutes
($d = 0.56$). French refereeing presents a distinctive institutional
feature: the Direction Technique de l'Arbitrage (DTA) manages
officiating for both divisions, and some referees work matches in both
competitions within a single season. This partial overlap in referee
pools means the enforcement signal did not need to cross an
organizational boundary to reach Ligue~2---it arrived through the same
officials who carried it to Ligue~1. The remaining 0.85-minute gap
likely reflects differences in match characteristics, VAR usage, and
broadcast scrutiny rather than institutional distance alone.

\paragraph{Italy: the failure.}
Serie~A added 0.73~minutes ($d = 0.40$); Serie~B showed no change
($\Delta = -0.08$, $d = -0.04$). The CAN managed Serie~A; a
separate body handled Serie~B. The muted top-flight response left
nothing to transmit downward.

\paragraph{Scotland: the media amplifier.}
The Scottish Premiership---nominally a control league---registered
the fourth-largest shift: $+1.47$~minutes ($d = 0.83$), outpacing
La~Liga. Scottish football has an independent referee pool, yet
operates in the same English-language media ecosystem as the Premier
League. Scottish assessors update their benchmarks watching Premier
League broadcasts; a Scottish referee adding only three minutes when
the Premier League routinely adds seven would feel conspicuously
inadequate. The Scottish case suggests the media channel may rival the
formal institutional one, though we cannot test this directly---we
lack, for example, evidence that Scottish assessor reports explicitly
reference Premier League benchmarks. The interpretation is suggestive
rather than established.

The within-country patterns converge: the enforcement gradient is
reproduced \textit{within} countries, between tiers sharing a
federation but differing in institutional proximity to the signal's
source.


% ──────────────────────────────────────────────────────────────
\subsection{The Italian Exception}
\label{sec:first_stage:italy}
% ──────────────────────────────────────────────────────────────

Serie~A's muted response demands sustained attention. By every formal
criterion---Big~Five membership, FIGC guidance, Italian referees at
the World Cup---it should sit near the gradient's summit. Instead, it
ranks eleventh of fourteen ($d = 0.40$), closer to the Greek Super
League ($d = 0.38$) and the Eredivisie ($d = 0.23$) than to any
Big~Five peer.

We do not have data to adjudicate among competing explanations, but
three features of Italian football are plausible contributors---each
speculative, offered to structure future inquiry rather than to
establish a mechanism. First, the \textit{calciopoli} legacy. The 2006
match-fixing scandal left Italian football acutely sensitive to
perceptions of referee discretion; adding substantially more stoppage
time is an inherently discretionary act that may carry professional
risks in this environment that are absent elsewhere. Second,
institutional insularity: the CAN operated at greater distance from
both the league and the international refereeing circuit than PGMOL
or the DFB referee department; its 2024 dissolution suggests Italian
football recognized the organizational limitations. Third, misaligned
media pressure: Italian discourse was consumed by VAR controversies
and referee-appointment politics rather than stoppage-time adequacy,
leaving the decentralized media-enforcement channel occupied by
other battles.

Match-level compliance sharpens the point: only 56.4\% of post-directive
Serie~A matches exceeded the league's pre-directive median---barely
above chance---versus 81.9\% in the Premier League. The directive
reached Serie~A but was absorbed rather than amplified. The Italian
exception confirms that the directive's force depended not on formal
designation but on the full apparatus through which enforcement signals
are translated into behavior on the pitch.


% ──────────────────────────────────────────────────────────────
\subsection{What Didn't Change}
\label{sec:first_stage:placebo}
% ──────────────────────────────────────────────────────────────

The specificity of the shift matters as much as its magnitude. Three
features of the data confirm that the change was confined to the
channel the directive activated.

\paragraph{First-half stoppage: the natural placebo.}
The directive applied to both halves, but cultural and media attention
focused on second-half stoppage---the added time that determines the
final whistle. If the shift reflected a mechanical trend in match
interruptions (more VAR reviews, longer injury assessments), first-half
stoppage would increase proportionally. It did not. First-half stoppage
showed a far smaller and statistically borderline change, consistent
with enforcement pressure concentrating where consequences of
non-compliance are most visible.

\paragraph{UEFA competitions: the institutional opt-out.}
UEFA did not adopt the new standard for the Champions League and
Europa League. The same referee who added nine minutes in a domestic
match on Saturday reverted to five in a continental fixture on
Tuesday---powerful evidence that the change was institutional, not
personal. The opt-out also eliminates the mechanical explanation:
Champions League matches use identical VAR technology, substitution
rules, and medical protocols. If the increase were driven by these
features rather than enforcement norms, it would appear in continental
fixtures. It did not.

\paragraph{Control-group contamination as evidence.}
Mean stoppage in the control group rose from 4.32 to 5.62~minutes, an
increase of 1.30~minutes---83\% of the treated-group shift. This is
not a nuisance for the story this section tells; it is the finding.
Thirteen of fourteen leagues shifted in the predicted direction. The
lone exception---Serie~B, at the bottom of a country whose top flight
itself barely moved---is the absence that confirms the pattern. The
consequences for identification are direct: a binary DiD yields a
conservative lower bound because the controls are themselves
substantially treated. Section~\ref{sec:strategy} develops the
within-country design and dose-response framework to accommodate this
reality.

\bigskip

The enforcement gradient documented here---from the Premier League's
2.11-minute shift at the summit to Serie~B's null response at the
base---is a story about how professional football is governed: who
controls referee behavior, through what channels enforcement signals
travel, and where the institutional architecture amplifies or dampens
the directives of the sport's global authority. The formal
identification strategy in Section~\ref{sec:strategy} exploits this
variation. But the institutional story comes first, because without it
the numbers are just numbers.
